\section{Ablation studies}
\label{sec:ablations}

\subsection{Compute--quality trade-off: MC passes and ensemble size}

Table~\ref{tab:compute_quality} varies the compute budget along two axes---MC
Dropout passes ($T$) and ensemble members ($N$)---to map the Pareto frontier
of uncertainty quality versus inference cost.

\begin{table}[H]
\centering
\caption{Compute--quality trade-off across MC Dropout passes and ensemble
sizes. Uncertainty quality plateaus after $T{=}10$; ensembles are
well-calibrated but slower than TTA at lower Unc-AUROC. TTA ($K{=}6$,
bottom row) dominates the frontier: highest Unc-AUROC at the lowest
multi-pass cost. Best value per column in bold.}
\label{tab:compute_quality}
\scriptsize
\setlength{\tabcolsep}{4pt}
\renewcommand{\arraystretch}{1.05}
\begin{center}
\begin{tabular}{l c c c c c}
\toprule
\textbf{Configuration} & \textbf{Dice} & \textbf{ECE} & \textbf{Unc-AUROC} &
\textbf{Time (s)} & \textbf{Unc/s} \\
\midrule
MC Drop.\ $T{=}5$   & 0.790 & 0.045 & 0.811 & \textbf{0.24} & 3.38 \\
MC Drop.\ $T{=}10$  & 0.792 & \textbf{0.042} & 0.836 & 0.47 & 1.78 \\
MC Drop.\ $T{=}20$  & 0.791 & 0.043 & 0.832 & 0.93 & 0.89 \\
MC Drop.\ $T{=}30$  & 0.791 & 0.044 & 0.838 & 1.42 & 0.59 \\
\addlinespace[2pt]
\midrule
\addlinespace[2pt]
Ensemble $N{=}2$     & 0.768 & 0.032 & 0.803 & 1.37 & 0.59 \\
Ensemble $N{=}5$     & \textbf{0.771} & \textbf{0.031} & 0.833 & 3.90 & 0.21 \\
\addlinespace[2pt]
\midrule
\addlinespace[2pt]
TTA $K{=}6$          & 0.768 & 0.035 & \textbf{0.881} & 0.86 & \textbf{1.02} \\
\bottomrule
\end{tabular}
\end{center}
\end{table}

The mean prediction stabilizes by $T = 5$: Dice varies by less than
0.002 across all configurations. Uncertainty quality benefits from more
passes, but with sharply diminishing returns. The jump from $T{=}5$ to
$T{=}10$ adds 2.5 points of Unc-AUROC; from $T{=}10$ to $T{=}30$, the
gain is only 0.2 points while runtime triples. For real-time applications,
$T{=}10$ is the practical sweet spot, retaining 99.8\% of the full
Unc-AUROC at one-third the cost.

Ensembles achieve the best ECE (0.031), confirming prior findings that
they are well-calibrated \citep{lakshminarayanan2017simple}, but this
calibration advantage does not translate to better deferral: a 5-member
ensemble (Unc-AUROC 0.833, 3.90\,s) is outperformed by TTA (Unc-AUROC
0.881, 0.86\,s) on both uncertainty quality and speed. The rightmost
column, Unc-AUROC per second, makes TTA's efficiency advantage explicit.

\subsection{Deferral threshold sensitivity}

We evaluate how sensitive each deferral strategy is to its threshold
parameter by measuring error reduction across a range of deferral rates.

\begin{table}[H]
\centering
\caption{Error reduction at fixed deferral budgets. Confidence-aware
deferral dominates at low review budgets (5--10\%), the regime most
relevant for clinical workflows, while adaptive thresholding wins at
higher budgets (20--30\%). Bold marks the best value per column.}
\label{tab:threshold_sensitivity}
\scriptsize
\setlength{\tabcolsep}{3.5pt}
\renewcommand{\arraystretch}{1.03}
\begin{center}
\begin{tabular}{l l c c c c}
\toprule
\textbf{Uncertainty} & \textbf{Deferral} &
\textbf{5\%} & \textbf{10\%} & \textbf{20\%} & \textbf{30\%} \\
\midrule
MC Dropout & Global     & 8.2\%  & 14.7\% & 23.5\% & 28.1\% \\
MC Dropout & Conf-aware & 21.4\% & 38.6\% & 52.1\% & 57.3\% \\
\addlinespace[2pt]
\midrule
\addlinespace[2pt]
TTA        & Global     & 15.3\% & 28.1\% & 44.8\% & 56.2\% \\
TTA        & Adaptive   & 18.7\% & 35.9\% & \textbf{61.3\%} & \textbf{75.4\%} \\
TTA        & Conf-aware & \textbf{24.8\%} & \textbf{43.2\%} & 60.7\% & 68.1\% \\
\bottomrule
\end{tabular}
\end{center}
\end{table}

At a 10\% deferral budget---a realistic clinical operating point---TTA with
confidence-aware deferral reduces errors by 43.2\%, nearly $3\times$ the
14.7\% achieved by MC Dropout with global thresholding. This three-fold
gap at the same review budget is the most practically relevant comparison
in the paper. Adaptive thresholding dominates at higher deferral rates
(20--30\%) because its per-image normalization ensures every image
contributes to the deferral pool; confidence-aware deferral dominates at
low rates because it concentrates the budget on the riskiest pixels.

\subsection{Cross-validation stability}
\label{sec:cv}

Five-fold cross-validation on the 20 DRIVE training images confirms that
segmentation metrics are highly stable: Dice $0.780 \pm 0.004$ (CV $<$ 1\%),
AUC $0.963 \pm 0.003$. ECE shows moderate variance ($0.044 \pm 0.006$,
CV 14\%), reflecting sensitivity to the specific train/val split.
Unc-AUROC has the highest variance ($0.768 \pm 0.057$, CV 7.4\%), ranging
from 0.71 to 0.83 across folds---driven by the interaction between the
learned decision boundary and the uncertainty threshold on a small
validation set. Crucially, the qualitative ordering (TTA $>$ MC Dropout;
confidence-aware $>$ global) holds across all five folds.
